\section{Описание}

Требуется реализовать алгоритм поиска лексикографически минимальной циклической ротации строки с использованием суффиксного дерева.

\subsection{Основная идея}

Суффиксное дерево - это сжатое дерево всех суффиксов строки. Для решения задачи о минимальной циклической ротации применяется классический приём: строится суффиксное дерево для строки $s + s + \$$ (где \$ --- non vocab символ), что позволяет рассмотреть все циклические сдвиги исходной строки как суффиксы этой конкатенации.

Для строки длины $n$ существует $n$ циклических ротаций. Наивный подход требует $O(n^2 \log n)$ времени для сравнения всех вариантов. Использование суффиксного дерева позволяет достичь линейной сложности $O(n)$ при поиске оптимальной ротации.

\subsection{Структура суффиксного дерева}

Каждый узел дерева содержит:
\begin{itemize}
    \item $l, r$ --- границы подстроки в исходной строке
    \item $\mathtt{next}$ --- словарь переходов по символам
    \item $\mathtt{slink}$ --- суффиксная ссылка для оптимизации построения
\end{itemize}

Суффиксная ссылка соединяет узел, соответствующий подстроке $c \alpha$, с узлом, соответствующим подстроке $\alpha$ (где $c$ --- символ, $\alpha$ --- подстрока).

\subsection{Построение суффиксного дерева (алгоритм Укконена)}

Алгоритм Укконена строит суффиксное дерево за $O(n)$ времени в два прохода:

\begin{enumerate}
    \item \textbf{Расширение:} На каждом шаге добавляется новый суффикс исходной строки
    \item \textbf{Оптимизация через суффиксные ссылки:} При несовпадении используются суффиксные ссылки вместо возврата в корень
\end{enumerate}

Ключевые параметры алгоритма:
\begin{itemize}
    \item $\mathtt{act\_v}$ --- активный узел
    \item $\mathtt{act\_pos}$ --- позиция на ребре активного узла
    \item $\mathtt{act\_len}$ --- длина активного пути
    \item $\mathtt{rem}$ --- количество оставшихся расширений
\end{itemize}

\subsection{Поиск минимальной ротации в дереве}

После построения суффиксного дерева для $s + s + \$$ выполняется поиск оптимального пути от корня:

\begin{enumerate}
    \item Начинаем с корня дерева
    \item На каждом шаге выбираем лексикографически первый переход (исключая \$)
    \item Берём из каждого ребра количество символов, необходимых для получения строки длины $|s|$
    \item Останавливаемся, когда накоплено ровно $|s|$ символов
\end{enumerate}

Это гарантирует лексикографическую минимальность, так как мы всегда выбираем наименьший доступный символ на каждом уровне.

\pagebreak

\section{Исходный код}

\subsection{Структура узла}

\begin{lstlisting}[language=C++]
struct Node {
    int l, r;
    int slink;
    map<char, int> next;

    Node(int l, int r) : l(l), r(r), slink(-1) {}
};
\end{lstlisting}

\subsection{Вспомогательная функция}

\begin{lstlisting}[language=C++]
int edge_length(int v) const {
    if (v == 0) return 0;

    int right = (nodes[v].r == -1) ? leaf_end : nodes[v].r;

    return right - nodes[v].l + 1;
}
\end{lstlisting}

Функция вычисляет длину ребра, ведущего в узел $v$. Для листовых узлов используется переменная $\mathtt{leaf\_end}$, которая динамически обновляется во время построения.

\subsection{Построение дерева}

Основной цикл построения суффиксного дерева:

\begin{lstlisting}[language=C++]
for (int i = 0; i < text.length(); ++i) {
    extend(i);
}
\end{lstlisting}

На каждом шаге $i$ к дереву добавляются все суффиксы, начинающиеся с позиции 0 и далее, которые заканчиваются в позиции $i$.

\subsection{Поиск оптимальной ротации}

\begin{lstlisting}[language=C++]
string optimal_lecs_split(int s_len) const {
    int curr = 0;
    int len = 0;
    string res;
    res.reserve(s_len);

    while (len < s_len) {
        int next_el = -1;
        
        for (auto const& [c, v_idx] : nodes[curr].next) {
            if (c == '$') continue; 
            next_el = v_idx;
            break;
        }

        if (next_el == -1) break; 

        int elen = edge_length(next_el);
        int take = min(elen, s_len - len);

        res += text.substr(nodes[next_el].l, take);
        len += take;
        curr = next_el;
    }
    return res;
}
\end{lstlisting}

Функция выбирает лексикографически минимальный переход на каждом узле и накапливает символы до достижения длины исходной строки.

\subsection{Главная функция решения}

\begin{lstlisting}[language=C++]
string solve(const string& s) {
    if (s.empty()) return "";

    string t = s + s + "$";
    SuffixTree st(t);

    return st.optimal_lecs_split(s.size());
}
\end{lstlisting}

Конкатенация $s + s$ обеспечивает наличие всех циклических ротаций как суффиксов новой строки. Завершающий символ \$ предотвращает нежелательные совпадения.

Сложность: Построение дерева --- $O(n)$, поиск минимальной ротации --- $O(n)$. Итоговая сложность: $O(n)$.

\pagebreak
